\chapter{Постановка задачи и критерии оценки}

\section{Климатические данные как объект анализа}

В контексте данной работы под климатическими данными понимаются наблюдательные и модельные данные, описывающие состояние атмосферы и связанных с ней поверхностных процессов. На практике такие данные включают ряды наблюдений метеостанций, продукты реанализа, архивы численных прогнозов и производные показатели, рассчитанные на регулярных пространственных сетках. С точки зрения машинного обучения этот класс данных сочетает свойства временных рядов, пространственных полей и многоуровневых структур.

Осадки являются одной из наиболее сложных переменных. Для них характерны высокая доля нулевых значений, резкие локальные всплески, сильная асимметрия распределения и зависимость от орографии, циркуляции и сезонности. Поэтому модель должна учитывать не только среднее отклонение прогноза, но и способность обнаруживать редкие опасные события. Для задач раннего предупреждения эта способность часто важнее, чем небольшое улучшение средней ошибки.

\section{Станционная и сеточная постановки}

В работе используются два уровня представления данных. В станционной постановке объект соответствует конкретной станции и моменту времени. Модель работает с вектором признаков: лагами осадков и температуры, статистиками по скользящим окнам, календарными и географическими признаками. Такой подход хорошо интерпретируем и вычислительно дешев, но пространственные связи между соседними районами учитываются только косвенно.

В сеточной постановке объектом является пространственное поле или фрагмент поля за несколько предшествующих моментов времени. Вход имеет тензорный вид \((T,C,H,W)\), где \(T\) — число лагов, \(C\) — число каналов, \(H,W\) — размеры области. Такая постановка требует моделей, способных выделять локальные пространственные шаблоны и более дальние зависимости. Поэтому здесь естественно использовать сверточные блоки, механизмы внимания и архитектуры типа U-Net/Transformer \cite{unet2015,vit2021,convnext2022}.

Станционный и сеточный уровни не конкурируют напрямую. Станционная модель полезна для локального прогноза по рядам наблюдений. Сеточная модель нужна, когда требуется пространственно согласованный прогноз и анализ ошибок на карте. В работе эти два уровня используются последовательно: сначала строится сильная табличная базовая модель, затем анализируется более сложная региональная нейросетевая постановка.

\section{Определение аномального события}

В работе аномальное осадочное событие определяется как превышение локального высококвантильного порога. Пусть \(y\) — целевая сумма осадков на заданном горизонте, а \(q_{0.9}\) — оценка 90-го процентиля исторического распределения для соответствующей станции или ячейки сетки. Тогда бинарная метка события задается как
\[
    z = \mathbb{I}\{y > q_{0.9}\}.
\]
Такое определение позволяет учитывать пространственную и климатическую неоднородность. Один и тот же абсолютный уровень осадков может быть обычным для влажной области и экстремальным для сухой области; квантильный порог учитывает это различие.

Из непрерывного прогноза \(\hat y\) можно получить вероятность события через сглаженное сравнение с порогом:
\[
    p = \sigma\left(\frac{\hat y - q_{0.9}}{T}\right),
\]
где \(\sigma\) — логистическая функция, а \(T>0\) — температурный параметр. Эта схема связывает регрессионный прогноз с задачей детекции редких событий без построения отдельного классификатора.

\section{Метрики качества}

Для непрерывного прогноза осадков могут использоваться RMSE, MAE и связанные с ними показатели. Однако в данной работе основной упор делается на событийную оценку, то есть на качество различения экстремальных и обычных случаев. Поэтому используются ROC-AUC, ROCSS, PR-AUC и Brier score \cite{wilks2011,gneiting2007,murphy1973brier,jolliffe2012forecast}.

ROC-AUC показывает, насколько хорошо модель ранжирует положительные и отрицательные случаи. ROC skill score выражается через ROC-AUC следующим образом:
\[
    \mathrm{ROCSS} = 2\cdot \mathrm{ROC\mbox{-}AUC} - 1.
\]
Значение ROCSS, близкое к нулю, соответствует случайному ранжированию, а положительное значение означает наличие полезного сигнала.

PR-AUC дополнительно важна для редких событий, поскольку она чувствительна к дисбалансу классов. Brier score оценивает среднеквадратичную ошибку вероятностного прогноза:
\[
    \mathrm{BS} = \frac{1}{N}\sum_{i=1}^{N}(p_i-z_i)^2.
\]
При анализе пространственных полей метрики дополнительно взвешиваются по площади ячеек через множитель, пропорциональный \(\cos\varphi\), где \(\varphi\) — широта. Это нужно потому, что ячейки широтно-долготной сетки имеют разную площадь \cite{weatherbench2}.

\section{Итоговая постановка экспериментов}

В первой экспериментальной части решается задача станционного прогноза суммарных осадков на горизонтах 24, 72, 120 и 240 часов. Прогноз строится по лаговым, агрегированным, календарным и географическим признакам. Результаты климатологии, линейной модели, LightGBM и XGBoost сравниваются по событийным метрикам.

Во второй экспериментальной части используется региональная нейросетевая модель для сеточных полей. Она принимает последовательность пространственных полей и выдает многогоризонтный прогноз осадков. Вероятность аномального события вычисляется относительно локальной карты порогов \(q_{0.9}\). Таким образом, обе части работы связаны общей формализацией аномалии, но используют разные представления данных и разные классы моделей.
